\chapter{Adverse Selection}\label{ch:b2:adverse-selection}

An annuity turns savings into an income paid for life. It is a good deal for someone who expects to live long
and a poor one for someone who does not, and buyers know more about their health than the insurer does. In the
United Kingdom, annuities were priced on age and sex only, but buyers could choose features: payments that rise
over time, or a guarantee that pays the estate if the buyer dies early. The chart shows how the buyers of each
feature died, compared with otherwise similar buyers of a plain annuity.

\begin{openingchart}{Change in the probability of dying within five years of buying an annuity, by feature of the
  annuity, percentage points, compared with buyers of a level, non-guaranteed annuity of the same age and sex
  (and, for the last bar, with women). One UK insurer, buyers up to 1993.}{Finkelstein and Poterba (2004), Table~2}
\begin{tikzpicture}
\begin{axis}[omrecord, xbar=0pt, bar width=6pt, width=0.7\linewidth, height=0.5\linewidth, xmin=-20,
    xmax=8, ymin=0.4, ymax=5.6, ymajorgrids=false, xlabel={Percentage points},
    ytick=data, yticklabels from table={figdata/bachelor-2/15-adverse-selection/out/annuities.csv}{feature},
    legend style={at={(0.5,-0.22)}, anchor=north, legend columns=2,
    /tikz/every even column/.append style={column sep=10pt}}, reverse legend]
  \addplot[fill=omThm!50, draw=omThm] table[x=voluntary, y=pos]
    {figdata/bachelor-2/15-adverse-selection/out/annuities.csv};
  \addplot[fill=omDef!55, draw=omDef] table[x=compulsory, y=pos]
    {figdata/bachelor-2/15-adverse-selection/out/annuities.csv};
  \draw[black!60] (axis cs:0,0.4) -- (axis cs:0,5.6);
  \legend{Voluntary market, Compulsory market}
\end{axis}
\end{tikzpicture}
\end{openingchart}

% ledger: b2:fp2004-t1, b2:fp2004-t2
% figdata: annuities_summary (escalating_vs_male 3.5, indexed_vs_male 4.2)
Buyers who chose payments that rise over time, in line with prices or at a fixed rate, were less likely to die
within five years: by 15 to 19 percentage points in the voluntary market. Buyers who chose a guarantee to their
estate were more likely to. In the voluntary market the effect of choosing rising payments was three and a half to
more than four times that of being a man, the one characteristic of health the insurer priced on. Buyers chose the product
that paid best given what they knew of their own health.

This is \omterm{def:b2:adverse-selection:adverse-selection}{adverse selection}: when one side of a market knows more about the quality or the risk of what is traded,
the price attracts the trades that are worst for the other side. This chapter shows how it can shrink or destroy a
market, how economists test for it, why it sometimes runs the other way, and what private arrangements and public
rules do about it, at what cost.

\begin{recall}
From the school volumes: hidden information. From \cref{ch:b2:expected-utility,ch:b2:insurance-and-risk-sharing}:
expected utility, the \omterm{def:b2:expected-utility:risk-premium}{risk premium}, and the demand for insurance.
\end{recall}

\begin{notation}[Symbols of this chapter]\label{not:b2:adverse-selection:symbols}
In the \omterm{def:b2:adverse-selection:lemons}{lemons model}, $\theta$ is the quality of a good, known to its seller. In insurance markets, $q$ is the share of people
insured, ranked from the highest willingness to pay; $D(q)$ is the willingness to pay of the $q$-th person, $\mathit{MC}(q)$ her
expected cost to the insurer and $\mathit{AC}(q)$ the average expected cost of the first $q$.
\end{notation}

\section{The lemons model}\label{sec:b2:adverse-selection:lemons}

\begin{definition}[Adverse selection]\label{def:b2:adverse-selection:adverse-selection}
\emph{Adverse selection}\index{adverse selection} occurs when one side of a trade has private information about the
quality or risk of what is traded, and at any given price the trades that take place are disproportionately those
that are worst for the uninformed side.
\end{definition}

\begin{definition}[Lemons model]\label{def:b2:adverse-selection:lemons}
The \emph{lemons model}\index{lemons model} is a market for a good of varying quality in which each seller knows the
quality of her good, buyers know only the distribution of quality, and every good sells at one price.
\end{definition}

Let quality $\theta$ be uniform on $[0, 1]$, a seller value her good at $\theta$ and a buyer at $1.2\,\theta$. With full information every
good would be sold, since each is worth more to a buyer than to its seller. At a single price $p \le 1$, only sellers with $\theta \le p$
sell, so the average quality on offer is $p/2$, which buyers value at $0.6\,p$.

\begin{proposition}[Collapse of the market]\label{prop:b2:adverse-selection:collapse}
In the \omterm{def:b2:adverse-selection:lemons}{lemons model} with quality uniform on $[0, 1]$, seller values $\theta$ and buyer values $1.2\,\theta$, no trade takes place at any
positive price.
\end{proposition}

\begin{proof}
At a price $p$ in $(0, 1]$, the goods offered have average quality $p/2$, worth $0.6\,p < p$ to a buyer, so no buyer pays $p$. At a price above
1 every good is offered, with average quality $\tfrac12$, worth $0.6 < 1 < p$. So at every positive price buyers do not buy, though every
trade would create value.
\end{proof}

\begin{omfigure}
\begin{tikzpicture}[econ]
\begin{axis}[omrecord, grid=none, xmin=0, xmax=1.05, ymin=0, ymax=1.05, width=0.5\linewidth,
    height=0.5\linewidth, xlabel={Price $p$}, ylabel={Value}, xtick={0,0.5,1}, ytick={0,0.5,1},
    legend style={at={(1.05,0.5)}, anchor=west}]
  \addplot[black, thick, domain=0:1, samples=2] {x};
  \addplot[omDef, thick, dashed, domain=0:1, samples=2] {x/2};
  \addplot[omThm, thick, domain=0:1, samples=2] {0.6*x};
  \legend{Price, Average quality offered, Buyers' value of it}
\end{axis}
\end{tikzpicture}

{\small The \omterm{def:b2:adverse-selection:lemons}{lemons model}. At a price $p$ the goods offered are those worth at most $p$ to their sellers, of average quality $p/2$.
Buyers value that average at $0.6\,p$, below the price at every positive price, so no price clears the market.}
\end{omfigure}

The good sellers leave first: at any price, the sellers of the best goods find it too low. The model does not say
that markets for used goods collapse; it says that, without ways for quality to be shown, they would shrink, and it
explains why sellers and buyers invest in such ways (\cref{sec:b2:adverse-selection:answers}).

\section{Insurance with hidden types}\label{sec:b2:adverse-selection:insurance}

Insurance is the case that matters most. Take two kinds of people, each facing a loss of 10{,}000 dalens. Low risks
lose it with probability 0.1 and would pay up to 1{,}500 a year for full cover; high risks lose it with probability 0.3 and
would pay up to 3{,}500. Each would pay more than her expected loss, the difference being her \omterm{def:b2:expected-utility:risk-premium}{risk premium}. There are as
many of each.

If the insurer could tell them apart it would charge 1{,}000 and 3{,}000, and everyone would buy. If it cannot, a single
premium that breaks even when everyone buys is 2{,}000. At that premium low risks leave, since 2{,}000 exceeds their 1{,}500; the
insurer, left with high risks only, must charge 3{,}000. The low risks go uninsured, losing a surplus of 500 each, though
insuring them costs less than they would pay.

\section{Unravelling}\label{sec:b2:adverse-selection:unravelling}

\begin{definition}[Unravelling]\label{def:b2:adverse-selection:unravelling}
\emph{Unravelling}\index{unravelling} (an adverse-selection spiral) is the process in which a rise in an insurance
premium drives out the buyers with the lowest expected costs, which raises the average cost of those who remain and
so the premium, driving out more.
\end{definition}

With many kinds of people, the market can be drawn with three curves. Rank people by willingness to pay for a
given policy, highest first: $D(q)$ is the demand curve. Let $\mathit{MC}(q)$ be the expected cost of the $q$-th person and $\mathit{AC}(q)$ the
average cost of the first $q$. Under \omterm{def:b2:adverse-selection:adverse-selection}{adverse selection} those who value cover most cost most, so $\mathit{MC}$ falls with $q$ and
$\mathit{AC}$ lies above it. Competitive insurers that cannot price on risk break even where demand meets average cost.
Coverage is efficient where demand meets marginal cost: a person should be insured when she values the cover more
than it costs to provide it to her.

Take $D(q) = 900 - 500q$ and $\mathit{MC}(q) = 800 - 600q$ (dalens a year), so that $\mathit{AC}(q) = 800 - 300q$. Every person values cover above her
cost, by $100 + 100q$, so everyone should be insured. But demand meets average cost where $900 - 500q = 800 - 300q$: half the people buy
at a premium of 650. The welfare lost is the area between demand and marginal cost over the uninsured half, $\int_{0.5}^1 (100 +
100q)\,dq = 87.5$ dalens a person.

\begin{omfigure}
\begin{tikzpicture}[econ]
\begin{axis}[omrecord, grid=none, xmin=0, xmax=1.05, ymin=150, ymax=950, width=0.62\linewidth,
    height=0.45\linewidth, xlabel={$q$: share insured}, ylabel={Dalens a year}, xtick={0,0.5,1},
    ytick={200,400,650,800}, legend style={at={(1.03,0.5)}, anchor=west}]
  \addplot[fill=omDef!15, draw=none, forget plot] coordinates {(0.5,650) (1,400) (1,200) (0.5,500)} -- cycle;
  \addplot[omDef, very thick, domain=0:1, samples=2] {900 - 500*x};
  \addplot[omThm, very thick, domain=0:1, samples=2] {800 - 300*x};
  \addplot[omMeth, very thick, dashed, domain=0:1, samples=2] {800 - 600*x};
  \node[eqpoint] at (axis cs:0.5,650) {};
  \draw[guide] (axis cs:0.5,150) -- (axis cs:0.5,650);
  \node[anchor=south west, font=\footnotesize] at (axis cs:0.5,655) {equilibrium};
  \node[font=\footnotesize] at (axis cs:0.8,400) {loss};
  \legend{Demand $D$, Average cost $\mathit{AC}$, Marginal cost $\mathit{MC}$}
\end{axis}
\end{tikzpicture}

{\small An insurance market with \omterm{def:b2:adverse-selection:adverse-selection}{adverse selection}. Insurers that cannot price on risk charge the average cost of their
buyers; the market settles where demand meets average cost, insuring half the people at 650. Everyone values cover above
her cost, so the shaded area between demand and marginal cost is the welfare lost.}
\end{omfigure}

How the market gets there shows the spiral. The listing starts an insurer at the premium that would break even if
everyone bought, and lets it reset the premium each year to the average cost of the last year's buyers.

\omcode{code/bachelor-2/15-adverse-selection/listings/unravelling.py}{1}{26}{Premiums and coverage year by year in a
market that unravels toward its equilibrium.}\label{lst:b2:adverse-selection:unravelling}

\omoutput{code/bachelor-2/15-adverse-selection/outputs/unravelling.txt}{The premium climbs and coverage falls, year by
year, to 650 and one half.}\label{out:b2:adverse-selection:unravelling}

% code: listings/unravelling.py
At 500 a fifth of the people find cover too dear; those who leave are the cheapest to insure, so the average cost of
those who stay rises to 560; more leave; and so on, each step smaller than the last. In this example the spiral stops at
half the market. If average cost lay above demand at every share but zero, it would not stop: the market would unravel
completely, as the lemons market does (\cref{exo:b2:adverse-selection:11}).

\begin{omfigure}
\begin{tikzpicture}[econ]
\begin{axis}[omrecord, grid=none, xmin=0.4, xmax=1.03, ymin=450, ymax=700, width=0.62\linewidth,
    height=0.42\linewidth, xlabel={$q$: share insured}, ylabel={Premium (dalens a year)},
    xtick={0.5,0.6,0.7,0.8,0.9,1}, ytick={500,560,596,650}]
  \addplot[omDef, very thick, domain=0.4:1, samples=2] {900 - 500*x};
  \addplot[omThm, very thick, domain=0.4:1, samples=2] {800 - 300*x};
  \addplot[black, thick, -{Stealth[length=2mm]}] coordinates {(1,500) (0.8,500) (0.8,560) (0.68,560) (0.68,596)
    (0.608,596) (0.608,617.6) (0.5648,617.6) (0.5648,630.56)};
  \node[eqpoint] at (axis cs:0.5,650) {};
  \node[anchor=east, font=\footnotesize, omDef] at (axis cs:0.56,560) {demand};
  \node[anchor=south, font=\footnotesize, omThm] at (axis cs:0.78,600) {average cost};
\end{axis}
\end{tikzpicture}

{\small \omterm{def:b2:adverse-selection:unravelling}{Unravelling} step by step. Each horizontal step is the buyers who leave at the current premium; each vertical step
is the rise in average cost when they have gone. The steps shrink toward the equilibrium at one half and 650.}
\end{omfigure}

\begin{model}{Demand and cost curves of an insurance market}
\assumes Insurers that cannot price on the risk that buyers know, competing until profits are zero; a fixed policy.
\predicts Coverage where demand meets average cost, below the efficient level when those who value cover most cost
most; a premium cut by a subsidy or mandate lowers average cost as low risks join.
\tested The Massachusetts individual mandate, after which premiums and average costs in the individual market
fell (Hackmann, Kolstad and Kowalski 2015).
\breaks When willingness to pay depends on other things than risk, such as aversion to risk or care, so that those who
value cover most need not cost most.
\end{model}

\section{Testing for adverse selection}\label{sec:b2:adverse-selection:testing}

\begin{definition}[Positive-correlation test]\label{def:b2:adverse-selection:positive-correlation}
The \emph{positive-correlation test}\index{positive-correlation test} checks whether, among people whom the insurer
treats alike, those who choose more coverage have more claims; asymmetric information predicts that they do.
\end{definition}

The test needs care. A positive correlation can also come from moral hazard, people with more cover taking less care
(\cref{ch:b2:moral-hazard}), and it must be computed among buyers whom the insurer prices alike, or it measures only
what the insurer already knows. A zero correlation is evidence against selection on risk, but not against private
information that bears on other things.

\begin{debate}[How much adverse selection is there in insurance?]
% ledger: b2:cs2012, b2:fp2004-t2
The answer differs by market. In car insurance, Chiappori and Salani\'e (2000) found, with three methods, no correlation
between coverage and claims in French car insurance, and most later work agreed; a positive correlation found in
Israel arose from drivers who, learning they were bad risks, switched insurers to buy underpriced cover (Cohen 2005). In
annuities, the choice of product is strongly related to later mortality (Finkelstein and Poterba 2004). In life insurance,
people who died within twelve years of a base year were 19 per cent more likely to have bought cover that year (He 2009). In
health insurance the results are mixed, and in one market the correlation is negative. All as surveyed by Chiappori and
Salani\'e (2012). One reading: \omterm{def:b2:adverse-selection:adverse-selection}{adverse selection} is real where insurers gather little information, as in annuities, and
small where they price finely, as in car insurance. Another: private information is everywhere, but it bears on several
traits at once, and they can offset each other.
\end{debate}

\section{Selection that runs the other way}\label{sec:b2:adverse-selection:advantageous}

\begin{definition}[Advantageous selection]\label{def:b2:adverse-selection:advantageous}
\emph{Advantageous selection}\index{advantageous selection} occurs when those who value insurance most have lower
expected costs, so that at a given premium the buyers are better risks than those who do not buy.
\end{definition}

People who are cautious, or averse to risk, may both buy more insurance and have fewer accidents. In the American market
for insurance that supplements public health cover for the elderly, people with higher cognitive ability were found to be
both more likely to buy and to have lower expected claims (Fang, Keane and Silverman 2008, as surveyed by Chiappori and
Salani\'e 2012). Then the marginal-cost curve rises with $q$ and the average cost lies below it. The market settles where
demand meets average cost, which is now to the right of where demand meets marginal cost: too \emph{many} people are
insured, some of whom value cover at less than it costs to provide it to them (\cref{exo:b2:adverse-selection:8}).
Selection on private information does not always shrink a market.

\section{Private and compulsory answers}\label{sec:b2:adverse-selection:answers}

Markets have their own ways to reduce the information gap.
\begin{itemize}
  \item \emph{Warranties and guarantees.} A seller who knows her good is sound can promise to repair it; a seller of a bad
    good cannot afford to.
  \item \emph{Certification and inspection.} Independent tests and ratings reveal quality before the sale.
  \item \emph{\omterm{def:b2:repeated-games:reputation}{Reputation}.} A dealer who sells many goods over many years has a record to lose (\cref{ch:b2:repeated-games}).
  \item \emph{Groups formed for other reasons.} Insurance sold to all employees of a firm covers people who joined the firm
    for reasons other than their risk, so the pool is not selected.
\end{itemize}

\begin{definition}[Risk classification]\label{def:b2:adverse-selection:risk-classification}
\emph{Risk classification}\index{risk classification} is the pricing of insurance by observable characteristics that
predict cost, such as age, sex, health history or driving record.
\end{definition}

\begin{definition}[Community rating]\label{def:b2:adverse-selection:community-rating}
\emph{Community rating}\index{community rating} is a rule requiring an insurer to charge the same premium for a
policy to all buyers, or to all in broad groups, whatever their observable risk.
\end{definition}

\begin{definition}[Insurance mandate]\label{def:b2:adverse-selection:mandate}
An \emph{insurance mandate}\index{insurance mandate} is a legal requirement to hold insurance, usually enforced by a
penalty on those who do not.
\end{definition}

Each answer has costs. \omterm{def:b2:adverse-selection:risk-classification}{Risk classification} narrows the information gap and with it the selection, but people with high
observed risk pay high premiums, a risk they could not insure before they were classified. \omterm{def:b2:adverse-selection:community-rating}{Community rating} protects them
but widens the gap, since the insurer may no longer use what it knows, and so strengthens selection. A mandate keeps low
risks in the pool and can restore full coverage; it also obliges people who value cover at less than its price to buy it,
and in a market with \omterm{def:b2:adverse-selection:advantageous}{advantageous selection}, or where some people value cover below its cost, it insures some who
should not be. Like any public rule, it depends on a penalty set and enforced by government, which may set it too low to
bind or too high.

\begin{record}[The Massachusetts mandate]
% ledger: b2:hkk2015
As part of a health reform, the state of Massachusetts required its residents to hold health insurance. Using data from the state, Hackmann, Kolstad and Kowalski (2015) estimated a model of selection and found that, in the
individual market, premiums and average costs fell significantly after the mandate, as healthier people joined. They
estimate an annual welfare gain from the reduction in \omterm{def:b2:adverse-selection:adverse-selection}{adverse selection} of 4.1 per cent per person, or 51.1 million dollars
a year in the state, and smaller markups after the reform.
\end{record}

\section{Exercises}\label{sec:b2:adverse-selection:exercises}

\begin{exercise}[$\star$]\label{exo:b2:adverse-selection:1}
In a lemons market quality is uniform between 0 and 10{,}000 dalens; a seller values her good at its quality and a buyer at 1.2
times it. At a price of 6{,}000, which sellers sell, what is the average quality on offer, and do buyers pay the price?
\end{exercise}

\begin{exercise}[$\star$]\label{exo:b2:adverse-selection:2}
Read the opening chart. Which features were chosen by buyers who lived longer, and which by buyers who died sooner? In which
market are the effects larger?
\end{exercise}

\begin{exercise}[$\star$]\label{exo:b2:adverse-selection:3}
Predict, then compute: in the two-type insurance market of the chapter, what is the smallest share of low risks at which a
single premium keeps both types insured?
\end{exercise}

\begin{exercise}[$\star$]\label{exo:b2:adverse-selection:4}
A study finds that drivers with more comprehensive cover have more accidents. List two explanations other than \omterm{def:b2:adverse-selection:adverse-selection}{adverse
selection}, and say what data would separate them.
\end{exercise}

\begin{exercise}[$\star\star$]\label{exo:b2:adverse-selection:5}
In the \omterm{def:b2:adverse-selection:lemons}{lemons model} let buyers value a good of quality $\theta$ at $b\theta$. For which $b$ is there trade at a positive price?
\end{exercise}

\begin{exercise}[$\star\star$]\label{exo:b2:adverse-selection:6}
In the market of \cref{sec:b2:adverse-selection:unravelling}, carry the spiral of the listing two more years by hand from a
premium of 596.
\end{exercise}

\begin{exercise}[$\star\star$]\label{exo:b2:adverse-selection:7}
Show that, for the curves $D(q) = 900 - 500q$ and $\mathit{MC}(q) = 800 - 600q$, the average cost is $800 - 300q$, and find the equilibrium and the
welfare loss.
\end{exercise}

\begin{exercise}[$\star\star$]\label{exo:b2:adverse-selection:8}
With \omterm{def:b2:adverse-selection:advantageous}{advantageous selection}, let $D(q) = 900 - 500q$ and $\mathit{MC}(q) = 300 + 200q$. Find the average cost, the equilibrium, the efficient
share insured and the welfare loss.
\end{exercise}

\begin{exercise}[$\star\star$]\label{exo:b2:adverse-selection:9}
Explain why the annuity buyers of the opening chart could select by product even though the insurer priced on age and sex.
Would pricing every feature on the mortality of its buyers end the selection?
\end{exercise}

\begin{exercise}[$\star\star\star$]\label{exo:b2:adverse-selection:10}
Prove that in the \omterm{def:b2:adverse-selection:lemons}{lemons model} with quality uniform on $[0, 1]$, seller values $\theta$ and buyer values $b\theta$ with $1 < b < 2$, no price
clears the market with positive trade, and find the gains from trade lost.
\end{exercise}

\begin{exercise}[$\star\star\star$]\label{exo:b2:adverse-selection:11}
With linear $D(q) = d_0 - d_1 q$ and $\mathit{MC}(q) = m_0 - m_1 q$, derive $\mathit{AC}(q)$ and the equilibrium share. Find the condition under which the
market unravels completely.
\end{exercise}

\begin{exercise}[$\star\star\star$]\label{exo:b2:adverse-selection:12}
A mandate imposes a penalty on people who stay uninsured. In the market of \cref{sec:b2:adverse-selection:unravelling}, find the
smallest penalty that makes everyone buy at the premium that breaks even when all are insured.
\end{exercise}

\section{Problem: Kestria's health cooperative}\label{sec:b2:adverse-selection:problem}

\begin{problem}[{Weekend problem --- a health cooperative in Kestria, the trading neighbour of Brenland, cannot tell its
members' health risks: how many stay insured, what does selection cost, and what do a mandate and risk-rated premiums
change}]
\label{pb:b2:adverse-selection:1}
The cooperative has 1{,}000 members and offers one policy, priced in kests a year to break even.

\medskip
\noindent\textbf{Part I --- Two risk types.} 600 members are healthy, with expected cost 400 and willingness to pay 600; 400 are
sick, with expected cost 1{,}600 and willingness to pay 1{,}900.
\begin{enumerate}
  \item What is the average expected cost of a member?
  \item What is the total surplus if everyone is insured?
  \item Should each type be insured, on grounds of efficiency?
  \item What premiums would the cooperative charge if it could tell the types apart?
  \item What is each type's \omterm{def:b2:expected-utility:risk-premium}{risk premium}?
  \item What information is hidden, and from whom?
\end{enumerate}

\medskip
\noindent\textbf{Part II --- A pooled premium.}
\begin{enumerate}[resume]
  \item At a premium of 880, who buys?
  \item What premium must the cooperative then charge, and who buys at it?
  \item How many members are insured?
  \item What surplus is lost?
  \item What surplus do the sick members keep?
  \item Is there any single premium at which everyone buys and the cooperative breaks even?
\end{enumerate}

\medskip
\noindent\textbf{Part III --- Demand and cost curves.} Now members differ continuously. Ranked by willingness to pay, the
member at rank $q$ would pay $1{,}800 - 1{,}000q$ and costs $1{,}500 - 1{,}200q$.
\begin{enumerate}[resume]
  \item Find the average cost of the first $q$ members.
  \item Find the equilibrium share insured and the premium.
  \item What is the efficient share insured?
  \item Find the welfare loss for the 1{,}000 members.
  \item Starting from a premium of 900, trace three years of the spiral.
  \item Why does the spiral stop short of \omterm{def:b2:adverse-selection:unravelling}{unravelling} completely?
\end{enumerate}

\medskip
\noindent\textbf{Part IV --- A mandate and risk-rated premiums.}
\begin{enumerate}[resume]
  \item Under a mandate everyone buys. What premium breaks even, and what is the gain in welfare?
  \item Which members pay more than they would choose to?
  \item What penalty makes the mandate bind for everyone?
  \item If the cooperative could observe each member's expected cost and charge it, who would buy, and what would the
    highest-cost member pay?
  \item Compare that member's premium under no intervention, the mandate and risk rating.
  \item State the result: equilibrium coverage and the welfare loss of selection, and the change under each answer.
\end{enumerate}
\end{problem}
